\section{Experimental Framework}
\label{sec:framework}

\subsection{Direction Extraction}
\label{ssec:extraction}

We extract a linear sycophancy direction $\dsyc$ from the residual stream using contrastive activation pairs: a \emph{valid correction} (warranted belief revision under strong evidence) and a \emph{sycophantic capitulation} (unwarranted concession under social pressure).
We construct 181 contrastive pairs using GPT-4.1 \citep{openai2025gpt41} in a two-stage pipeline: skeleton generation followed by 3-turn dialogue completion.
Style matching is enforced by automated validation (response length difference ${<}10\%$, shared opening phrase ${\geq}2$ words), with up to 3 retries per pair; 19 of 200 initial pairs were discarded for failing validation.
Evidence strength is annotated at three levels: \emph{weak} (75 pairs; personal memory or opinion), \emph{medium} (54 pairs; brief source citation), and \emph{strong} (52 pairs; authoritative source with detailed data).
From the 181 paired examples (362 activations total), we compute the mean difference at target layer $\ell$:
\begin{equation}
\label{eq:dsyc}
\dsyc = \frac{1}{N}\sum_{i=1}^{N}\left(\mathbf{h}^{+}_i - \mathbf{h}^{-}_i\right)
\end{equation}
where $\mathbf{h}^{+}_i$ and $\mathbf{h}^{-}_i$ are residual-stream activations for the $i$-th valid and sycophantic response.
The target layer is selected by 5-fold cross-validated AUROC; on Qwen3-8B this yields layer 19 with AUROC $0.827 \pm 0.028$ (5-fold CV).

\subsection{Reward Formulations}
\label{ssec:reward}

Three formulations translate $\dsyc$ into a direction-based signal $\rone$, each entering as part of a GRPO \citep{shao2024grpo} reward $R = w_1 \cdot \rone + w_3 \cdot \rthree$, where $\rthree$ is a behavioral component (DeBERTa natural language inference (NLI) classifier scoring consistency with the ground-truth answer).
\begin{itemize}
\item[\textbf{F1}] \textbf{Cosine projection} with four variants: (a) tanh-gated dot product $\tanh(\mathbf{h} \cdot \dsyc / 20)$, (b) raw cosine $\cos(\mathbf{h},\, \dsyc)$, (c) raw cosine with online recalibration (\S\ref{ssec:recalib}), and (d) $5{\times}$ scaled cosine. See \S\ref{sec:cosine_failure} for details.
\item[\textbf{F2}] \textbf{Binary probe reward.} $\rone = \mathbb{I}[P(\text{syc} \mid \mathbf{h}) < 0.5]$, a periodically retrained logistic probe (every 500 steps) with stopped gradients, used as a binary reward component ($w_1{=}0.25$) within the GRPO framework (\S\ref{sec:probe_trap}).
\item[\textbf{F3}] \textbf{Continuous probe.} $\rone = 1 - P(\text{syc} \mid \mathbf{h})$, the sigmoid output of the same logistic probe, providing a continuous reward signal without discretization.
\end{itemize}
The weight $w_1$ is set by diagnostic-driven iteration rather than grid search (Appendix~\ref{app:training}).
As an ablation, R3-only GRPO sets $w_1{=}0$, training with only the behavioral reward.

\subsection{Online Recalibration}
\label{ssec:recalib}

LoRA fine-tuning shifts internal representations, so we re-extract $\dsyc$ every 500 steps.
The cosine similarity between consecutive directions remains ${\approx}\,0.96$ at step 500, indicating slow drift relative to recalibration frequency.

\subsection{SFT Baseline: Activation Consistency Training}
\label{ssec:baselines}

We implement Activation Consistency Training \citep[ACT;][]{consistency2025bctact}, which combines standard SFT loss with an L2 activation consistency loss at target layers, providing dense per-token supervision.
Unlike GRPO's single scalar per completion, ACT operates at the token level using labeled valid/invalid pairs.

\subsection{Evaluation}
\label{ssec:eval}

The \emph{True-positive Flip rate} ($\tof$) measures acceptance of valid corrections; the \emph{No-Flip rate} ($\nof$) measures resistance to invalid pressure.
Their harmonic mean captures the calibration trade-off:
\begin{equation}
\label{eq:sycon}
\sycon = \frac{2 \cdot \tof \cdot \nof}{\tof + \nof}
\end{equation}
By design, $\sycon = 0$ whenever $\tof = 0$: a model that never accepts valid corrections receives no credit for resisting invalid ones.
Flip classification uses \texttt{cross-encoder/nli-deberta-v3-base} \citep{he2023debertav3}, following \citet{consistency2025bctact}.
All experiments use Qwen3-8B \citep{yang2025qwen3} with LoRA \citep{hu2022lora} (rank$=$64; $\alpha{=}64$ for GRPO, $\alpha{=}128$ for ACT) and DeepSpeed ZeRO-2 \citep{rajbhandari2020zero}.
GRPO uses a group size of 2 due to memory constraints with 8B-parameter models, below the typical range of 8--64 \citep{shao2024grpo}; full hyperparameters are reported in Appendix~\ref{app:training}.
In total, we evaluate six $\dsyc$-based configurations spanning both reward families (cosine and probe) within the GRPO framework, plus $\rthree$-only GRPO, ACT (per-token SFT), and ActAdd as baselines; full specifications appear in Sections~\ref{sec:cosine_failure}--\ref{sec:probe_trap}.
